\chapter{Números complejos}

\section{Cuerpos}

\begin{definicion}
Sea $K$ un conjunto y sean $+,\cdot : K\times K\to K$ dos operaciones en $K$.
Se dice que $(K,+,\cdot)$ es un \textbf{cuerpo} si se cumple:
\begin{enumerate}
    \item $+$ y $\cdot$ son asociativas y conmutativas.
    \item Existe un elemento neutro para la suma, $0\in K$, tal que $x+0=x$
    para todo $x\in K$.
    \item Existe un elemento neutro para el producto, $1\in K$, con $1\neq 0$,
    tal que $x\cdot 1=x$ para todo $x\in K$.
    \item Cualquiera sea $x\in K$, existe un inverso aditivo u \textbf{opuesto}
    $-x\in K$ tal que $x+(-x)=0$.
    \item Cualquiera sea $x\in K$ con $x\neq 0$, existe un inverso
    multiplicativo $x^{-1}\in K$ tal que $x\cdot x^{-1}=1$.
    \item El producto es distributivo sobre la suma:
    $x\cdot(a+b)=x\cdot a+x\cdot b$ para todos $x,a,b\in K$.
\end{enumerate}
\end{definicion}

\begin{ejemplo}
$\mathbb{Q}$, $\mathbb{R}$ y $\mathbb{C}$ son cuerpos con la suma y el producto
usuales. $\mathbb{Z}$ no lo es: por ejemplo $2\in\mathbb{Z}$ no tiene inverso
multiplicativo en $\mathbb{Z}$.
\end{ejemplo}

\section{Números complejos}

\subsection{Forma binomial}

\begin{definicion}
El conjunto de los \textbf{números complejos} es
\[
\mathbb{C}=\{a+bi : a,b\in\mathbb{R}\},
\]
donde $i$ es un símbolo que cumple $i^2=-1$, con la suma y el producto
usuales (distribuyendo y usando $i^2=-1$).
\end{definicion}

\begin{definicion}
Dado $z\in\mathbb{C}$, la expresión $z=a+bi$ con $a,b\in\mathbb{R}$ se llama la
\textbf{forma binomial} de $z$. Su \textbf{parte real} es $\operatorname{Re}(z)=a$
y su \textbf{parte imaginaria} es $\operatorname{Im}(z)=b$.

Dado $z=a+bi$ con $a,b\in\mathbb{R}$:
\begin{itemize}
    \item el \textbf{conjugado} de $z$ es $\overline{z}=a-bi\in\mathbb{C}$;
    \item el \textbf{módulo} de $z$ es $|z|=\sqrt{a^2+b^2}\in\mathbb{R}_{\ge 0}$.
\end{itemize}
\end{definicion}

\begin{proposicion}
Sea $z\in\mathbb{C}$. Entonces:
\begin{enumerate}
    \item $z\cdot\overline{z}=|z|^2$.
    \item Si $z\neq 0$, entonces $z^{-1}=\dfrac{\overline{z}}{|z|^2}$.
\end{enumerate}
\end{proposicion}

\begin{demostracion}
Sea $z=a+bi$ con $a,b\in\mathbb{R}$.

(1) $z\overline{z}=(a+bi)(a-bi)=a^2-abi+abi-b^2i^2=a^2+b^2=|z|^2$, usando $i^2=-1$.

(2) Si $z\neq 0$ entonces $|z|^2=a^2+b^2\neq 0$, y por (1)
\[
z\cdot\frac{\overline{z}}{|z|^2}=\frac{z\overline{z}}{|z|^2}=\frac{|z|^2}{|z|^2}=1,
\]
así que $\overline{z}/|z|^2$ es el inverso multiplicativo de $z$.
\end{demostracion}

\begin{ejemplo}
Calcular $\dfrac{1}{3+4i}$. Como $|3+4i|^2=9+16=25$,
\[
(3+4i)^{-1}=\frac{3-4i}{25}=\frac{3}{25}-\frac{4}{25}i.
\]
\end{ejemplo}

\begin{proposicion}[{title={Propiedades del conjugado y del módulo}}]
Sean $z,w\in\mathbb{C}$. Entonces:

\textbf{Conjugado.}
\begin{enumerate}
    \item $\overline{z+w}=\overline{z}+\overline{w}$.
    \item $\overline{zw}=\overline{z}\,\overline{w}$. En particular,
    $\overline{z^n}=\overline{z}^{\,n}$ para todo $n\in\mathbb{N}$.
    \item $\overline{\overline{z}}=z$.
    \item $z=\overline{z}\iff z\in\mathbb{R}$.
    \item $\operatorname{Re}(z)=\dfrac{z+\overline{z}}{2}$ y
    $\operatorname{Im}(z)=\dfrac{z-\overline{z}}{2i}$.
\end{enumerate}

\textbf{Módulo.}
\begin{enumerate}
    \setcounter{enumi}{5}
    \item $|z|\ge 0$, y $|z|=0\iff z=0$.
    \item $|\overline{z}|=|z|$.
    \item $|zw|=|z|\,|w|$. En particular, si $z\neq 0$, $|z^{-1}|=\dfrac{1}{|z|}$.
    \item $|\operatorname{Re}(z)|\le|z|$ y $|\operatorname{Im}(z)|\le|z|$.
    \item (Desigualdad triangular) $|z+w|\le|z|+|w|$.
\end{enumerate}
\end{proposicion}



\begin{proposicion}[{title={Parte real e imaginaria del inverso}}]
Sea $z=a+bi\neq 0$ con $a,b\in\mathbb{R}$. Entonces
\[
z^{-1}=\frac{\overline{z}}{|z|^2}=\frac{a-bi}{a^2+b^2},
\]
y por lo tanto
\[
\operatorname{Re}(z^{-1})=\frac{a}{a^2+b^2},
\qquad
\operatorname{Im}(z^{-1})=-\frac{b}{a^2+b^2}.
\]
\end{proposicion}

\begin{demostracion}
Es la fórmula $z^{-1}=\overline{z}/|z|^2$ vista antes (multiplicar arriba y abajo
por $\overline{z}$ para sacar la $i$ del denominador), escribiendo
$\overline{z}=a-bi$ y $|z|^2=a^2+b^2$ y separando parte real e imaginaria.
\end{demostracion}

\begin{proposicion}[{title={Multiplicar por $i$}}]
Sea $z=a+bi$ con $a,b\in\mathbb{R}$. Entonces
\[
i\cdot z=-b+ai,
\]
y por lo tanto
\[
\operatorname{Re}(iz)=-\operatorname{Im}(z),
\qquad
\operatorname{Im}(iz)=\operatorname{Re}(z).
\]
\end{proposicion}

\begin{demostracion}
$i\cdot z=i(a+bi)=ai+bi^2=-b+ai$, usando $i^2=-1$. Como $-b$ y $a$ son reales,
$\operatorname{Re}(iz)=-b=-\operatorname{Im}(z)$ e $\operatorname{Im}(iz)=a=\operatorname{Re}(z)$.
\end{demostracion}

\begin{nota}{}
Geométricamente, multiplicar por $i$ gira al punto $(a,b)$ un ángulo de
$90^\circ$ alrededor del origen, y lo manda a $(-b,a)$.
\end{nota}

\begin{proposicion}[{title={Potencias de $i$}}]
Las potencias de $i$ se repiten en un ciclo de $4$ resultados:
\[
i^1=i,\qquad i^2=-1,\qquad i^3=-i,\qquad i^4=1.
\]
Más aún, para todo $n\in\mathbb{N}_0$,
\[
i^n=i^{r_4(n)},
\]
donde $r_4(n)\in\{0,1,2,3\}$ es el resto de dividir $n$ por $4$.
\end{proposicion}

\begin{demostracion}
$i^2=-1$ por definición, $i^3=i^2\cdot i=-i$ e $i^4=i^2\cdot i^2=(-1)(-1)=1$.

Por el algoritmo de la división, $n=4q+r$ con $r=r_4(n)$ y $0\le r<4$. Entonces
\[
i^n=i^{4q+r}=(i^4)^q\, i^r=1^q\, i^r=i^r.
\]
\end{demostracion}

\begin{ejemplo}
Calcular $i^{2026}$. Como $2026=4\cdot 506+2$, se tiene $r_4(2026)=2$ y
\[
i^{2026}=i^2=-1.
\]
\end{ejemplo}


 
\subsection{Forma trigonométrica}
 
\begin{nota}{Periodicidad.}
Las funciones seno y coseno son periódicas de período
$2\pi$: al sumarle a un ángulo $\theta$ cualquier cantidad de vueltas completas
$2k\pi$ (con $k\in\mathbb{Z}$), el valor no cambia:
\[
\cos\theta=\cos(\theta+2k\pi), \qquad
\operatorname{sen}\theta=\operatorname{sen}(\theta+2k\pi), \qquad \forall k\in\mathbb{Z}.
\]
\end{nota}
 
\begin{definicion}
Sea $z\in\mathbb{C}$, $z\neq 0$. Una escritura de la forma
\[
z=r(\cos\theta+i\operatorname{sen}\theta), \qquad r\in\mathbb{R}_{>0},\ \theta\in\mathbb{R},
\]
se llama una \textbf{forma trigonométrica} (o \textbf{polar}) de $z$. Aquí $r$ es
el \textbf{módulo} de $z$ (la distancia al origen, por eso es estrictamente
positivo) y $\theta$ es un \textbf{argumento} de $z$ (un ángulo en radianes).
\end{definicion}
 
\begin{proposicion}
Sea $z=a+bi\in\mathbb{C}$ con $z\neq 0$. Entonces $z$ admite una forma
trigonométrica $z=r(\cos\theta+i\operatorname{sen}\theta)$, y necesariamente
$r=|z|$.
\end{proposicion}
 
\begin{demostracion}
\textbf{El módulo.} Si $z=r(\cos\theta+i\operatorname{sen}\theta)$ con $r>0$, entonces
\[
|z|=r\sqrt{\cos^2\theta+\operatorname{sen}^2\theta}=r.
\]
 
\textbf{Existencia.} Tomamos $r=|z|>0$. Como
$\left(\tfrac{a}{r}\right)^2+\left(\tfrac{b}{r}\right)^2=\tfrac{a^2+b^2}{r^2}=1$,
el punto $\left(\tfrac{a}{r},\tfrac{b}{r}\right)$ está en la circunferencia
unitaria, así que existe $\theta\in\mathbb{R}$ con $\cos\theta=\tfrac{a}{r}$ y
$\operatorname{sen}\theta=\tfrac{b}{r}$. Entonces
$r(\cos\theta+i\operatorname{sen}\theta)=a+bi=z$.
\end{demostracion}
 
\begin{proposicion}[{title={Igualdad en forma trigonométrica}}]
Sean $z,w\in\mathbb{C}$ no nulos con
\[
z=r(\cos\theta+i\operatorname{sen}\theta), \qquad
w=s(\cos\varphi+i\operatorname{sen}\varphi),
\]
donde $r,s\in\mathbb{R}_{>0}$ y $\theta,\varphi\in\mathbb{R}$. Entonces
\[
w=z \iff
\begin{cases}
r=s\\
\varphi=\theta+2k\pi \ \text{ para algún } k\in\mathbb{Z}.
\end{cases}
\]
\end{proposicion}
 
\begin{demostracion}
($\Leftarrow$) Si $r=s$ y $\varphi=\theta+2k\pi$, por la periodicidad
$\cos\varphi=\cos\theta$ y $\operatorname{sen}\varphi=\operatorname{sen}\theta$,
así que $w=z$.
 
($\Rightarrow$) Si $w=z$, entonces $|w|=|z|$, es decir $s=r$. Cancelando $r>0$ en
$r(\cos\theta+i\operatorname{sen}\theta)=r(\cos\varphi+i\operatorname{sen}\varphi)$
e igualando partes real e imaginaria queda
\[
\cos\varphi=\cos\theta, \qquad \operatorname{sen}\varphi=\operatorname{sen}\theta.
\]
Entonces
\[
\cos(\varphi-\theta)=\cos\varphi\cos\theta+\operatorname{sen}\varphi\operatorname{sen}\theta
=\cos^2\theta+\operatorname{sen}^2\theta=1,
\]
y como $\cos x=1$ si y sólo si $x=2k\pi$ con $k\in\mathbb{Z}$, resulta
$\varphi=\theta+2k\pi$ para algún $k\in\mathbb{Z}$.
\end{demostracion}
 
\begin{definicion}
Por la proposición anterior, los argumentos de un $z\neq 0$ están determinados
salvo múltiplos de $2\pi$, y hay exactamente uno con $0\le\theta<2\pi$. Ese
único valor se llama el \textbf{argumento} de $z$ y se anota $\arg(z)$.
\end{definicion}
 
\begin{nota}{Valores notables.}
\begin{center}
\renewcommand{\arraystretch}{1.8}
\begin{tabular}{|c|c|c|c|c|c|c|}
\hline
$\theta$ & $0$ & $\frac{\pi}{6}$ & $\frac{\pi}{4}$ & $\frac{\pi}{3}$ & $\frac{\pi}{2}$ & $\pi$ \\
\hline
$\cos\theta$ & $1$ & $\frac{\sqrt{3}}{2}$ & $\frac{\sqrt{2}}{2}$ & $\frac{1}{2}$ & $0$ & $-1$ \\
\hline
$\operatorname{sen}\theta$ & $0$ & $\frac{1}{2}$ & $\frac{\sqrt{2}}{2}$ & $\frac{\sqrt{3}}{2}$ & $1$ & $0$ \\
\hline
\end{tabular}
\end{center}
\end{nota}
 
\begin{ejemplo}
Escribir en forma trigonométrica $z=1+i$ y $w=1+\sqrt{3}\,i$.
 
Para $z$: $|z|=\sqrt{1+1}=\sqrt{2}$, y buscamos $\theta$ con
$\cos\theta=\tfrac{1}{\sqrt{2}}=\tfrac{\sqrt{2}}{2}$ y
$\operatorname{sen}\theta=\tfrac{\sqrt{2}}{2}$, que es $\theta=\tfrac{\pi}{4}$. Luego
\[
1+i=\sqrt{2}\left(\cos\tfrac{\pi}{4}+i\operatorname{sen}\tfrac{\pi}{4}\right),
\qquad \arg(z)=\tfrac{\pi}{4}.
\]
 
Para $w$: $|w|=\sqrt{1+3}=2$, con $\cos\theta=\tfrac{1}{2}$ y
$\operatorname{sen}\theta=\tfrac{\sqrt{3}}{2}$, que es $\theta=\tfrac{\pi}{3}$. Luego
\[
1+\sqrt{3}\,i=2\left(\cos\tfrac{\pi}{3}+i\operatorname{sen}\tfrac{\pi}{3}\right),
\qquad \arg(w)=\tfrac{\pi}{3}.
\]
\end{ejemplo}
 
\begin{proposicion}[{title={Conjugado e inverso en forma trigonométrica}}]
Sea $z=r(\cos\theta+i\operatorname{sen}\theta)$ con $r>0$ y $\theta\in\mathbb{R}$.
Entonces:
\begin{enumerate}
    \item $\overline{z}=r\bigl(\cos(-\theta)+i\operatorname{sen}(-\theta)\bigr)$.
    \item $z^{-1}=\dfrac{1}{r}\bigl(\cos(-\theta)+i\operatorname{sen}(-\theta)\bigr)$.
\end{enumerate}
\end{proposicion}
 
\begin{demostracion}
(1) Por definición de conjugado,
\[
\overline{z}=r(\cos\theta-i\operatorname{sen}\theta)
=r\bigl(\cos(-\theta)+i\operatorname{sen}(-\theta)\bigr),
\]
pues el coseno es par ($\cos(-\theta)=\cos\theta$) y el seno es impar
($\operatorname{sen}(-\theta)=-\operatorname{sen}\theta$).
 
(2) Como $|z|=r$, por la fórmula $z^{-1}=\overline{z}/|z|^2$ y por (1),
\[
z^{-1}=\frac{\overline{z}}{|z|^2}
=\frac{r\bigl(\cos(-\theta)+i\operatorname{sen}(-\theta)\bigr)}{r^2}
=\frac{1}{r}\bigl(\cos(-\theta)+i\operatorname{sen}(-\theta)\bigr).
\]
\end{demostracion}
 
\begin{nota}{}
Geométricamente, $\overline{z}$ es el simétrico de $z$ respecto del eje real: tiene
el mismo módulo $r$ y el argumento opuesto $-\theta$. Si $\theta\in(0,2\pi)$,
el argumento de $\overline{z}$ en $[0,2\pi)$ es $2\pi-\theta$. Y $z^{-1}$ tiene
módulo $1/r$ y también argumento $-\theta$.
\end{nota}
 